\chapter{Model selection without the truth}
\label{ch:model-selection}

Every estimator in the paper has a knob that is not a parameter of the model: the
number of EM iterations, or the number of gradient steps, at which the fit is read off.
The error of a fit is a function of that knob, usually falling and then rising, and the
value reported for an arm is its error at the knob's chosen position. How the position
is chosen is therefore part of the estimator, and the paper spends a surprising amount
of its protocol on it: validation and test bundles, checkpoint ladders, a selection
rule per arm, a screening stage on seeds that never enter a reported number, and a
control experiment (E20) whose only purpose is to show that one arm's selection does
not depend on the truth. This chapter defines each of these objects, explains why they
are there, and gives the three information interfaces A, B and C that the paper uses to
say what each learner was allowed to see.

\section{Bundles: what a fit sees and what it is scored on}
\label{sec:bundles}

\subsection{Three disjoint samples}

A \term{bundle} is a fixed set of sequences together with everything derived from them
that an experiment needs. The paper keeps three per cell:
\begin{itemize}
\item the \term{training bundle}, $\nseq$ clean chains $\va^1,\dots,\va^{\nseq}$ drawn from
  the prior, from which each arm builds whatever it is allowed to see (\Cref{sec:interfaces});
\item the \term{validation bundle}, $1024$ chains, used only to choose the knob;
\item the \term{test bundle}, $2048$ chains, used only to score the chosen fit.
\end{itemize}
For the validation and test bundles the code draws the chains once and then, for every
level $t$ of the twelve-level schedule, one noised view $\vx=e^{-t}\va+\sqrt{\Dt}\vz$ of
every chain and the reference posterior mean $\pmean{\va}$ of that view under the
\emph{true} kernel, computed by the tree-exact recursion of \Cref{ch:factor-graphs-bp}
on the frozen grid (\code{\_bundle} in \code{experiments/e03\_baselines.py}; the same
function shape in E02, E08, E14). A bundle is therefore a dictionary
$t\mapsto(\mat{X}_t,\mat{M}_t)$ of the noised views and their reference means at each
level, and its SHA-256 hash over those two arrays is stored with every row of every CSV,
so that two arms claiming to be scored on the same bundle can be checked to have been
(\Cref{ch:protocol}).

\subsection{Keyed random streams}

The three bundles are drawn from three random streams keyed by a tag, the seed and, for
the training bundle, the size. In the code the training stream is
\code{rng\_for("exp31-train", seed, M)}, the validation stream is
\code{rng\_for("exp31-val", seed)} and the test stream is
\code{rng\_for("exp31-test", seed)}. Two consequences
matter for model selection. First, the streams are disjoint by construction (different
tags), so no validation or test chain is ever a training chain. Second, they are the
same across experiments: E02a's EM--BP fit and E03's network fit at $(\text{seed},\nseq)$
see the identical training chains and are scored on the identical test bundle, which is
what makes their ratio a \term{paired} comparison rather than a difference of two
averages (\Cref{ch:experimental-statistics}). The registry calls the pair of tags and
sizes the \emph{bundle protocol}; \Cref{ch:protocol} explains the tag mechanism.

\subsection{The one-view groups}

EM--BP never sees the clean training chains. From them the code derives the
\term{one-view groups}: a random permutation of the $\nseq$ chains is split into twelve
consecutive blocks of (nearly) equal size, block $k$ is assigned to level $t_k$, and
each chain in it receives \emph{one} noised view at that level, drawn once from the
stream \code{rng\_for("exp31-em-noise", seed)} and never redrawn
(\code{\_one\_view\_groups}). The result is a list of twelve pairs
$(\mat{X}_{t_k},t_k)$ with about $\nseq/12$ rows each. This is the data of interface~C.
A held-out one-view bundle for selection is built the same way from the validation
chains with a different noise stream.

\section{Checkpoint ladders}
\label{sec:ladders}

\begin{definition}[Checkpoint ladder]
\label{def:ladder}
A \term{checkpoint ladder} is a fixed, increasing list of budgets at which a fit is
saved. For EM--BP it is the list of iteration counts
\begin{equation}
  \{10,\,20,\,40,\,60,\,80,\,120,\,160,\,220,\,300,\,400,\,600,\,800,\,1200\},
\end{equation}
and for the networks the list of gradient-step counts
\begin{equation}
  \{500,\,1000,\,2000,\,4000,\,8000,\,16000,\,32000\}.
\end{equation}
A \term{rung} is one entry; the
\term{cap} is the last rung. The \term{selected checkpoint} is the rung the selection
rule picks, and a cell is \term{at cap} when the selected rung is the last one.
\end{definition}

Two things the definition does not say deserve saying. The ladder is coarse on purpose:
its rungs are spaced roughly geometrically, so the selection can tell a factor of two
in budget apart but not ten per cent. And the ladder is not a stopping rule: the fit is
always run to the cap (unless EM's own likelihood tolerance ends it first,
\Cref{ch:maximum-likelihood-em}), the kernel or network state at every rung is saved, and
the choice is made afterwards, with every rung available. This is what fixed the
stopping-rule artefact of the earlier capacity study, where one arm stopped at its
tolerance near its optimum and the others ran to a fixed cap past theirs (see the paper,
Appendix on the previous capacity reading).

The at-cap fraction of a set of cells is a diagnostic the paper reports whenever it
bears on a claim: if the selected rung is the cap in most cells, validation was still
improving when the ladder ran out, and the reported error is an upper bound on what a
longer ladder would give. The paper says so for the largest networks at
$\nseq=2048$ and for EM--BP at the two highest levels of the learning experiment; the
budget controls E02a\_v1 and E06c exist to settle whether the cap binds.

\section{Oracle selection: validation risk against the true kernel}
\label{sec:oracle-selection}

\begin{definition}[Oracle selection rule]
\label{def:oracle-rule}
Let $\hat m^{(r)}(\vx,t)$ be the posterior-mean map of the fit at rung $r$. Its
\term{validation risk} is the schedule-pooled relative score error of
\Cref{ch:experimental-statistics},
\begin{equation}
  R_{\mathrm{val}}(r)=\Biggl(\frac{\sum_t\sum_{j}\sum_{i\in B}\bigl(\hat s^{(r)}_t(\vx_j)_i-s_t(\vx_j)_i\bigr)^2}{\sum_t\sum_j\sum_{i\in B}s_t(\vx_j)_i^2}\Biggr)^{1/2},
  \label{eq:val-risk}
\end{equation}
where $s_t(\vx)=-(\vx-e^{-t}\pmean{\va})/\Dt$ and $\hat s^{(r)}_t$ is the same expression
with $\hat m^{(r)}$ in place of $\pmean{\va}$,
all computed on the validation bundle over the \term{bulk} region $B$ of sites $5$ to
$\nsites-4$, with $s_t$ the reference score built from the true kernel's posterior means
stored in the bundle. The \term{oracle rule} selects $r^\ast=\argmin_rR_{\mathrm{val}}(r)$.
\end{definition}

The rule is called an oracle because $s_t$ is the score of the true law, which a
practitioner does not have. Three remarks.

\paragraph{Why the bulk.} Every arm is scored on all sites for the headline, but
selected on the interior. The boundary sites are where the arms differ for reasons
unrelated to what is being selected (a zero-padded convolution behaves differently at
the ends; the chain's first site has no left neighbour), and a selection statistic that
weights them would let a boundary artefact choose the rung. The paper's protocol
appendix records the three regions (all, bulk, centre) and which one is used for what.

\paragraph{Why it is legitimate for a comparison.} The oracle rule is applied
identically to every arm: EM--BP at every capacity, the six architectures under
interface~A, the fixed-corruption arm. A comparison between arms that all selected on
the same oracle statistic is fair as a comparison; what it cannot support is a claim
that any single arm could be deployed at its reported error without the truth. The
paper draws exactly this line: the only arm selected without an oracle is the
implicit-score-matching arm, and E20 checks EM--BP.

\paragraph{Selection is a shrinkage.} Choosing the rung on validation is early
stopping, and early stopping is a regulariser: an EM run stopped before convergence
sits between the initialisation and the maximum-likelihood estimate. The audit's
per-level reading of the learning experiment (\Cref{ch:em-through-noise}) turns on this:
at high noise the selected fits are shrunk toward the Gaussian initialisation, which is
why they can sit \emph{below} an unbiased-estimator bound.

\section{Oracle-free selection for EM--BP: held-out log-evidence}
\label{sec:evidence-selection}

\subsection{The statistic}

\begin{definition}[Held-out one-view log-evidence per edge]
\label{def:evidence}
Let $\{(\mat{X}_{t_k},t_k)\}$ be a held-out one-view bundle of $n$ chains and let
$\kernelth^{(r)}$ be the kernel at rung $r$. The \term{held-out log-evidence per edge} is
\begin{equation}
  E(r)=\frac{1}{n(\nsites-1)}\sum_{k}\sum_{\vx\in\mat{X}_{t_k}}\log P_{t_k,\params^{(r)}}(\vx),
  \label{eq:evidence}
\end{equation}
the average over held-out views of the marginal log-likelihood of the view under the
fitted kernel, divided by the number of transitions per chain. The \term{evidence rule}
selects $r^\ast=\argmax_rE(r)$.
\end{definition}

The marginal log-likelihood $\log P_{t,\params}(\vx)$ is exactly what the forward sweep
of the discretised model returns when its normalising constants are accumulated
(\Cref{ch:grid-quadrature}; the field \code{log\_evidence} returned by
\code{e\_step\_multi} in \code{src/markov\_score/em.py}). No clean chain and no true kernel enters: the statistic
is a function of the fitted kernel and of noised data only. Dividing by the number of
edges makes it comparable across chain lengths and sizes (E14 uses it at $\nsites=128$).

\subsection{Why maximising held-out likelihood is a valid rule}

For a held-out sample from the true noised law $\Pt$, the average log-likelihood of a
candidate law $Q$ is an unbiased estimate of $\E_{\Pt}\log Q(\vx)$, and
\begin{equation}
  \E_{\Pt}\log Q(\vx)=\E_{\Pt}\log\Pt(\vx)-\KL(\Pt\,\|\,Q).
  \label{eq:heldout-kl}
\end{equation}
The first term does not depend on $Q$. So the rung with the largest held-out evidence
is, up to the sampling error of the held-out set, the rung whose fitted law is closest
to the truth in Kullback--Leibler divergence. It is the same principle as
cross-validation of a likelihood, and it is oracle-free because \cref{eq:heldout-kl}
never requires $\Pt$ to be known, only sampled.

Two caveats keep the rule honest. The criterion is Kullback--Leibler divergence of the
noised law, not the relative score error the paper reports; the two are different
functionals of $Q$ and can rank two rungs differently. And the held-out set is finite,
so $E(r)$ carries sampling noise of order $n^{-1/2}$ per edge, which can blur the choice
between neighbouring rungs whose evidence differs by less than that. Both are visible in
the control.

\subsection{E20}

E20 (\code{experiments/e20\_evidence\_selection.py}) replicates E02a cells exactly, at
$\ncomp=8$ with the same bundle tags, the same initialisation stream and the same ladder,
and records at every rung three numbers: the oracle validation risk
\cref{eq:val-risk}, the held-out one-view evidence \cref{eq:evidence} computed on the
validation chains with a fresh noise stream (tags \code{e20-split}, \code{e20-noise}),
and the test risk at the rungs the two rules and the cap pick. The comparison is the
test risk under the evidence rule against the test risk under the oracle rule, paired
by seed, and the fraction of seeds in which the two rules pick the same rung. The design
concentrates on $\nseq=32$, the size at which the choice of rung matters most (at
$\nseq=2048$ the frozen validation curves are monotone through the ladder, so every rule
picks the cap and the control is empty); the $\nseq=512$ arm was dropped for cost, as
the registry records. The paper quotes the result in Section~5 and Appendix A5: the two
rules pick the same rung in most seeds, give the same test error within seed noise,
and both are far below the error at the cap. The sentence it licenses is that EM--BP's
reported numbers do not depend on the oracle, not that the oracle rule was never used.

\section{Oracle-free selection for networks: the implicit objective}
\label{sec:ism-selection}

\subsection{Hyv\"arinen's identity}

A network trained by implicit score matching (\Cref{ch:score-matching,ch:neural-baselines})
minimises, over noisy samples $\vx\sim\Pt$ only,
\begin{equation}
  J(\theta)=\E_{\Pt}\Bigl[\tfrac12\lVert s_\theta(\vx)\rVert^2+\nabla\!\cdot s_\theta(\vx)\Bigr].
  \label{eq:ism}
\end{equation}
The reason this is a score-matching objective is an integration by parts. For a
smooth density $p$ and a smooth vector field $s_\theta$ whose product with $p$ vanishes
at infinity,
\begin{equation}
  \E_p[\nabla\!\cdot s_\theta]=\int(\nabla\!\cdot s_\theta)\,p\dd\vx=-\int s_\theta\cdot\nabla p\dd\vx
  =-\E_p[s_\theta\cdot\nabla\log p],
\end{equation}
so
\begin{equation}
  J(\theta)=\tfrac12\E_p\lVert s_\theta\rVert^2-\E_p[s_\theta\cdot s_p]
  =\tfrac12\E_p\lVert s_\theta-s_p\rVert^2-\tfrac12\E_p\lVert s_p\rVert^2 ,
  \label{eq:hyvarinen}
\end{equation}
with $s_p=\nabla\log p$ the true score. The last term does not depend on $\theta$.
Minimising $J$ is minimising the \term{Fisher divergence} $\tfrac12\E_p\lVert s_\theta-s_p\rVert^2$
without ever evaluating $s_p$.

\subsection{As a selection rule}

The same identity makes $J$ a model-selection criterion. Evaluate
\cref{eq:ism} on a held-out one-view bundle: the empirical average is an unbiased
estimate of $J$ at that rung, hence of the Fisher divergence up to a constant common to
all rungs. Selecting $r^\ast=\argmin_rJ_{\mathrm{held}}(r)$ picks the rung closest to the
true score in Fisher divergence, with no clean chain and no true score anywhere
(\code{ism\_objective} in \code{src/markov\_score/nn\_torch.py}). Three details of the
paper's implementation:
\begin{itemize}
\item the divergence in the selection statistic is computed \emph{exactly}: by the
  coloured probes for the window head and by $\nsites$ backward passes for a dense map
  (\Cref{ch:neural-baselines}); a stochastic estimate would add noise to a comparison
  between rungs that may differ by little;
\item the held-out bundle carries one view per chain split across the twelve levels,
  exactly as the training data, so the criterion weights the levels as the data do:
  equal numbers of sequences per level, each contributing its own
  $\tfrac12\E\lVert s_\theta-s_p\rVert^2$; the paper's estimand instead pools squared
  errors and divides by the pooled reference energy, which weights the low levels more
  (\Cref{ch:experimental-statistics}), so the two can rank rungs differently;
\item the whole held-out curve $r\mapsto J_{\mathrm{held}}(r)$ is stored in the CSV, not
  only its minimum, so that a run whose curve never turns over can be recognised as a
  divergence of the optimisation rather than reported as the network's error; E03 saw
  this at $\nseq=32$, where implicit score matching on a tiny sample is known to be
  unstable, and the paper says so.
\end{itemize}

\section{Screening and confirmation}
\label{sec:screening}

\subsection{Two kinds of seeds}

The paper's protocol names two disjoint sets of seeds. The \term{development seeds}
$10000$ to $10003$ (E03's screen used the first two after a pre-data amendment) are for
every choice a designer makes: architecture family, hyperparameters, learning rate,
output parameterisation. The \term{confirmatory seeds} $0$ to $15$ (eight of them in the
regime sweeps) are for the reported numbers. The rule is stated twice in the contract
and once in the paper: development seeds never enter a reported number; confirmatory
seeds never enter a selection.

\subsection{Why: the winner's curse}

Suppose $K$ configurations are tried on the same data and the best one's error is
reported. Even if all $K$ have the same true error, the reported minimum is biased
downward, because it is the minimum of $K$ noisy measurements. The paper avoids this by
choosing on one set of seeds and reporting on another: the confirmatory error of the
chosen configuration is then an ordinary, unbiased estimate of that configuration's
error, because the choice did not look at the confirmatory data.

\begin{example}[The winner's curse for two candidates, by hand]
\label{ex:winners-curse}
Two configurations have the same true risk $R$, and their measured risks on the
development seeds are $R+\sigma Z_1$ and $R+\sigma Z_2$ with $Z_1,Z_2$ independent
standard normal. The screen reports $\min(R+\sigma Z_1,R+\sigma Z_2)=R-\sigma\max(-Z_1,-Z_2)$.
For two independent standard normals, $\max(Z_1,Z_2)$ has expectation
$\E|Z_1-Z_2|/2=\tfrac12\cdot2/\sqrt{\pi}\cdot\tfrac{1}{\sqrt2}\cdot\sqrt2=1/\sqrt\pi$: indeed
$\max(a,b)=\tfrac12(a+b)+\tfrac12|a-b|$, the first term has mean zero, and $Z_1-Z_2$ is
normal with variance $2$, so $\E|Z_1-Z_2|=\sqrt{2}\cdot\sqrt{2/\pi}=2/\sqrt\pi$. Hence the
screened minimum has expectation $R-\sigma/\sqrt\pi\approx R-0.56\sigma$: with a
screen noise of, say, five per cent of the risk, reporting the screen's own number
would understate the risk by nearly three per cent for two candidates, and by more for
the sixty-eight configurations of E03's screen. Reporting the winner's risk on fresh
seeds removes the bias entirely; the fresh measurement has its own noise but no
preference for the small side.
\end{example}

\subsection{What a fair screen is}

\begin{definition}[Fair screen]
\label{def:fair-screen}
A screen is \term{fair} when every candidate is trained under the same recipe (loss,
optimiser, batch, gradient clipping, level sampling, corruption), for the same budget
(steps and restarts), on the same bundles, with the same checkpoint ladder and the same
selection statistic, and the winner of each family is chosen by the same rule on
development seeds disjoint from the confirmatory ones.
\end{definition}

The definition is a list because each item names a way the earlier screen of this
model was unfair, as the audit found: the long-range architectures had been trained with
one noise level per gradient step where the window head had twelve, and with one
restart where the head had eight, so a difference between families was partly a
difference in training. E03's screen (\code{part\_screen} in
\code{experiments/e03\_baselines.py}) gives every family the recipe of
\Cref{ch:neural-baselines}: per-example level from the twelve, batch $64$, Adam,
clipping at $1$, both parameterisations, learning rates $\{10^{-3},3\cdot10^{-3}\}$, a
family-specific hyperparameter grid (radius and width for the window head; channels for
the TCN; hidden size for the GRU; width, heads and depth for the transformer; channels
for the U-Net; hidden widths for the MLP), one run of $16{,}000$ steps per configuration
at the two sizes $\nseq\in\{32,2048\}$ on two development seeds, the checkpoint selected
by the oracle rule on the bulk, and the family's winner chosen by the pooled bulk risk
over development seeds and sizes (\code{winners\_from\_screening}). The winners are
written to \code{outputs/E03/winners.json} and are the configurations of every
confirmatory row. The paper reports the screen in its protocol appendix (the screen
table) and says that the screen's numbers are never a headline.

\subsection{Confirmation}

The \term{confirmatory run} (\code{part\_confirm}) trains each family's winner on the
sixteen confirmatory seeds at the four sizes $\{32,128,512,2048\}$ for the full
$32{,}000$ steps, selects the rung by the oracle rule, and scores on the test bundle in
all three regions. Its rows are what the baseline table shows and what the paired
ratios against EM--BP are computed from, joined on the seed, the size and the test-bundle
hash.

\section{The three information interfaces}
\label{sec:interfaces}

The paper's comparisons are organised by what a learner is allowed to see. The
definitions below are the code's, made explicit.

\begin{definition}[Interfaces A, B and C]
\label{def:interfaces}
Let $\{\va^\mu\}_{\mu\le\nseq}$ be the training chains and let each chain be assigned one
level $t_\mu$ by the even split of \Cref{sec:bundles}.
\begin{itemize}
\item \term{Interface A} (the standard denoising recipe). The learner may query the clean
  chains at every gradient step, together with a freshly drawn level $t$ (uniform on the
  twelve) and a freshly drawn corruption $\vz$; its training pairs are
  $(\va^\mu,\ e^{-t}\va^\mu+\sqrt{\Dt}\vz)$ with new $(t,\vz)$ every time. The checkpoint
  is selected by the oracle rule.
\item \term{Interface B} (one fixed corruption). The learner sees the pairs
  $(\va^\mu,\vx^\mu)$ with $\vx^\mu=e^{-t_\mu}\va^\mu+\sqrt{\Delta_{t_\mu}}\vz^\mu$, one
  corruption per chain, drawn once. It still regresses on the clean chain. The
  checkpoint is selected by the oracle rule.
\item \term{Interface C} (one noised view per sequence and nothing else). The learner
  sees $\{(\vx^\mu,t_\mu)\}$ only: the one-view groups. It never sees a clean chain, and
  its checkpoint is selected without the true score, by the held-out implicit
  objective (networks) or the held-out evidence (EM--BP, E20 and E14's second row).
\end{itemize}
\end{definition}

The three are nested in what they reveal: C sees a function of what B sees, and B sees
one draw of what A can draw indefinitely. The paper's Section~5 uses the nesting in one
sentence: the gap between the interface-A rows and the interface-C row is the value of
clean data with fresh corruptions; the factor that remains at C is what the
factorisation and the exact E-step are worth at equal information. Interface~B is the
control that separates the two ingredients of A: it gives the clean target but not the
augmentation, so a network at B that is still far behind EM--BP shows that one
corruption per sequence, not the absence of the clean chain, is what it cannot recover
from. The audit's correction of 18 September was precisely a relabelling: an arm that
sees $\va^\mu$ is at B, not at C, whatever its loss is called, and the headline
``matched'' number has to be the interface-C arm.

EM--BP is at interface~C by construction: its E-step reads $\{(\vx^\mu,t_\mu)\}$ and
nothing else (\Cref{ch:em-through-noise}). Its selection rule in the capacity and
baseline experiments is the oracle one, which is why E20 exists; in E14 both rules are
run side by side.

\section{The regime sweep of the matched comparison: E14}
\label{sec:e14}

E14 (\code{experiments/e14\_regime\_baselines.py}) repeats the interface-C comparison
away from the frozen point: correlations $\corr\in\{0.5,0.85,0.95\}$, chain lengths
$\nsites\in\{32,128\}$, Laplace innovations, $\nseq=2048$, eight confirmatory seeds, with
its own bundles (tags \code{e14-train}, \code{e14-val}, \code{e14-test},
\code{e14-valchains}, keyed by seed and length). Each cell scores six arms on a common
test bundle:
\begin{description}[leftmargin=2.2em,itemsep=0.25em]
\item[\code{em\_bp}] the eight-component mixture kernel fitted by EM--BP on the one-view
  groups, rung selected by the held-out BP risk of the validation bundle against the
  true kernel's posterior means (the oracle rule of \Cref{def:oracle-rule}, on all sites
  here);
\item[\code{em\_bp\_evidence}] the same EM run and the same ladder, rung selected by the
  held-out one-view log-evidence per edge on the noised views of the validation chains
  (the evidence rule of \Cref{def:evidence}): the oracle-free EM--BP;
\item[\code{ism\_one\_view}] a network trained by implicit score matching on the same
  one-view groups, rung selected on the held-out implicit objective with the exact
  divergence; one row per architecture family (window head, transformer), each with the
  hyperparameters that won a screen \emph{inside the regime} on two development seeds
  over E03's candidate grid, $8{,}000$ steps for the screen and $32{,}000$ for the
  confirmation; the window head trains with the coloured-probe divergence and learning
  rate $10^{-3}$, the transformer with Hutchinson's estimator and $3\cdot10^{-3}$, as
  E03's winners did;
\item[\code{dsm\_one\_corruption}] the same architectures at interface~B: denoising score
  matching on clean targets paired with the very same one noised view per chain
  (the index arrays of the split are returned so the pairing is exact), rung selected by
  the validation BP risk;
\item[\code{lmmse\_true}] the LMMSE estimator with the true covariance, the second-order
  floor (\Cref{ch:gaussian-closure-lmmse});
\item[\code{isotropic}] the score $-\vx$ of the standard Gaussian, the reference of a
  model that learned nothing beyond the unit variance.
\end{description}
The predeclared analysis is, per cell, the mean pooled relative score error of every arm
and the paired ratio of the best network to EM--BP over seeds with a percentile
bootstrap, with the trend in $\corr$ reported against the marginal's excess kurtosis
$3(1-\corr^2)/(1+\corr^2)$ as a direction, not a fitted law. The registry entry carries a
pre-data amendment worth reading as a specimen of the protocol: the first screen had
trained the window head with Hutchinson's one-probe divergence and obtained risks far
above the same network's E03 value; the cause was identified as the estimator, not the
regime, the screen was kept as a negative result, and the design was re-run with the
coloured probes before any confirmatory cell had run (\Cref{ch:protocol}). At the time
of writing E14 is queued on the cluster and no number from it is in the paper; the
chapter describes the design, not a result.

\where{\begingroup\sloppy The paper, Section~5 (the three interfaces; the oracle-free
selection sentence), Appendix A3 (arms and interfaces; the screen table) and Appendix A5
(E20); \Cref{ch:neural-baselines} for the losses and \Cref{ch:em-through-noise} for the
EM side. Code: \code{experiments/protocol.py} (bundle tags, ladders, regions, dev and
confirmatory seeds), \code{experiments/e03\_baselines.py} (screen, confirm, matched,
length), \code{experiments/e20\_evidence\_selection.py},
\code{experiments/e14\_regime\_baselines.py}, \code{experiments/e02\_capacity.py},
\code{src/markov\_score/nn\_torch.py} (\code{ism\_objective}),
\code{src/markov\_score/em.py} (\code{e\_step\_multi}, \code{fit\_em}),
\code{tools/make\_baseline\_numbers.py} (the matched macros by arm).\par\endgroup}

\paragraph{References.}
A.~Hyv\"arinen, 2005, Estimation of non-normalized statistical models by score matching, \emph{Journal of Machine Learning Research} 6, 695--709.
M.~Stone, 1974, Cross-validatory choice and assessment of statistical predictions, \emph{Journal of the Royal Statistical Society B} 36, 111--147.
L.~Prechelt, 1998, Early stopping, but when?, in \emph{Neural Networks: Tricks of the Trade}, Springer, 55--69.
G.~C.~Cawley and N.~L.~C.~Talbot, 2010, On over-fitting in model selection and subsequent selection bias in performance evaluation, \emph{Journal of Machine Learning Research} 11, 2079--2107.
P.~Vincent, 2011, A connection between score matching and denoising autoencoders, \emph{Neural Computation} 23, 1661--1674.
